\section{Introduction}
\label{sec:introduction}

A sampled position stream can look accurate at its endpoints while being dynamically wrong between them. Consider a constant-speed reference sampled every control period. A follower sends the next position to a jerk-limited online trajectory generator (OTG), executes the first period of the resulting trajectory, and repeats. If every target silently includes zero terminal velocity and acceleration, the generator is not being asked to continue the motion. It is being asked to stop at every sample. When that stop is reachable within one period, the exact within-period solution is a short rest-to-rest pulse: accelerate, decelerate, arrive at the correct position, and reach zero velocity before the next command. Repetition produces \stopandgo{} motion even though the endpoint position error can be nearly zero. This is not feedback chattering, stick--slip, or a failure of jerk limitation. It is a mismatch between the intended continuous motion and the terminal derivatives carried by the command.\claimtag{C1}

This distinction matters in streaming robot interfaces. OTG algorithms are normally formulated for complete initial and target states under kinematic bounds. Ruckig, for example, accepts target position, velocity, and acceleration and computes jerk-limited trajectories to arbitrary target states~\cite{berscheid2021ruckig}. An application that supplies positions while leaving the derivative fields at zero has therefore made a substantive modeling choice, whether or not the API makes that choice visually prominent. The resulting artifact may be missed by a benchmark based only on target-time position root-mean-square error (RMSE). It is visible in the exact piecewise-jerk profile, in the within-period velocity minimum and ripple, and in the repeated terminal stops.

The phenomenon raises three questions that are not answered by the instruction to ``set target velocity.'' First, when is a one-period rest-to-rest move reachable under acceleration and jerk limits? Without that boundary, the occurrence of the artifact appears parameter-specific. Second, does a matched velocity target remove the mechanism because it carries the correct state semantics, or merely because it provides another tunable input? Third, can the velocity be constructed causally from a scheduled position stream, and over what noise and timing envelope does that construction remain useful? These questions lead to a theory--experiment separation that is central to this paper. The theory concerns a single third-order integrator and proves reachability or feasible continuation. The experiments characterize how one pinned implementation, Ruckig~\RuckigVersion, realizes those possibilities. No theorem asserts that every OTG uses the same profile-selection policy.

We derive the maximum displacement of a symmetric jerk- and acceleration-limited rest-to-rest transfer in one control period. With period $T$, bounds $A$ and $J$, and reference speed $\vref$, the result reduces to two dimensionless variables,
\begin{equation}
 q=\frac{4A}{JT}, \qquad
 \rho=\frac{\abs{\vref}}{\vcrit}.
 \label{eq:intro-dimensionless}
\end{equation}
The boundary $\rho=1$ separates one-period rest-to-rest reachability from nonreachability when the velocity bound is inactive. The formula also gives a counterintuitive design implication: increasing $A$, $J$, or $T$ weakly enlarges the set of speeds for which repeated stopping is reachable. Looser dynamic limits can therefore expose, rather than cure, a bad terminal-state contract.\claimtag{C2}

The structural remedy is equally simple to state and more precise than a tuning heuristic. If the current state is $(p_k,\vref,0)$ and the next target is $(p_k+\vref T,\vref,0)$, then constant velocity with zero acceleration and jerk is an exact feasible continuation. We call this a \matchedvelocity. Its importance is not that velocity is always the best extra field, but that the terminal state is consistent with the invariant motion the application intends. We further specify a fixed-grid information contract in which $P[k+1]$ is an already scheduled target, while measurements are available only through index $k$. Under that contract, \FutureOOne{} constructs a target velocity from $P[k],P[k-1],P[k-2]$ without reading a future measurement.\claimtag{C3}\claimtag{C5}

The empirical program is organized by hypotheses rather than by experiment chronology. A boundary study completes \EFifteenGridCount{} required cells and \EFifteenSobolCount{} out-of-grid threshold searches; it confirms the dimensionless transition away from an explicitly reported set of \EFifteenSeamCount{} exact-seam numerical failures. A causal intervention study executes \ESixteenArmCount{} arms, varying the velocity coefficient and comparing wrong-sign, random-sign, position-lookahead, and minimum-duration alternatives. Only the matched velocity interventions reproduce the preregistered exact profile in every tested primary cell. A frozen robustness study then evaluates a development-selected causal observer on \ESeventeenWorkPairCount{} paired seed--configuration cells across 11 declared conditions. The weakest declared condition still has a median ripple reduction of \ESeventeenWeakMedianPercent, while six stronger, out-of-envelope diagnostics expose a failure boundary rather than being hidden. Finally, one recorded fixed-grid trajectory provides an engineering case study: PV \FutureOOne{} reduces RMSE by \RecordedPVReductionPercent{} and the \observedlag{} from 20 to \RecordedPVLag. A timestamp-aware local-polynomial replay does not improve the irregular-timestamp version, and that negative result remains part of the main account.\claimtag{C4}\claimtag{C6}\claimtag{C9}\claimtag{C13}

The contributions are fourfold:
\begin{enumerate}[leftmargin=*,itemsep=2pt]
  \item a closed-form one-period rest-to-rest reachability boundary, its $q$--$\rho$ normalization, and monotonic design consequences under explicit assumptions;
  \item a proof that the matched PV contract admits an invariant zero-jerk continuation, separated from empirical claims about a particular solver;
  \item exact-profile causal ablations and frozen holdout tests that reject several competing explanations and state both a tested work envelope and its observed boundary; and
  \item a reproducible recorded case study that distinguishes waveform tracking, observed lag, numerical guardrails, and best-tested constraint selection.
\end{enumerate}

The scope is deliberately narrower than a deployment claim. The theory is single-axis. The experiments are software evaluations without plant dynamics or multi-axis synchronization. The recorded evidence contains one waveform, and measured offline-host execution time is not a real-time schedulability result. These limits let the central causal claim remain testable: a position endpoint can be correct while its derivative contract makes the continuous motion wrong.

The remainder of the paper reviews OTG, reference shaping, and causal derivative estimation; defines the receding information contract and metrics; develops the dimensionless analysis; presents the boundary, ablation, robustness, and recorded results; and closes with limitations and implications for streaming motion interfaces.

